\subsection*{13.2 Expectation Conditioned on a Sub-sigma-algebra}

We extend the notion of conditional expectation to the general case by replacing the

special\sigma-algebra. \mathcal{G}in Definition 13.2 by any sub-\sigma-algebra of. \mathcal{F}.

\begin{defbox}{13.3}
\end{defbox}

Let X be a real-valued P -integrable random variable and \mathcal{G}is a sub-\sigma-algebra

of \mathcal{F}The conditional expectation of X given \mathcal{G}, denoted by E[X\vert\mathcal{G}],i sa \mathcal{G}-

measurable random variable that satisfies

\int

B

E[X\vert\mathcal{G}] dP =

\int

B

XdP for all B \in\mathcal{G}. (13.3)

The condition in (13.3) can be written in two equivalent forms. These forms

provide alternative ways to think about the definition of E[X\vert\mathcal{G}].

(i)

\int

\Omega

E[X\vert\mathcal{G}]\cdot 1B dP =

\int

\Omega

X \cdot 1B dP for all B \in\mathcal{G}; (13.4)

(ii)

\int

\Omega

E[X\vert\mathcal{G}]\cdot ZdP =

\int

\Omega

X\cdot ZdP (13.5)

for all bounded \mathcal{G}-measurable random variables Z .

It is obvious that condition (ii) in (13.5) implies condition (i) in (13.4) because

the indicator function . 1B for any B\in \mathcal{G} is \mathcal{G}-measurable. To show the reverse

direction, we can proceed with the sequence of indicator functions, simple functions,

and bounded functions that are \mathcal{G}-measurable.

Comparing condition (ii) in (13.5) with the version of Orthogonal Principle

in (12.5), we see that when X is square integrable, the conditional expectation of

X given a sub-\sigma-algebra \mathcal{G}is the same as the projection of X on the space of \mathcal{G}-

measurable functions with finite second moments. Roughly speaking, for random

variables in L2(P) , the conditional expectation is just a minimum mean-squared

error (MMSE) estimator. A more precise statement is given in the next theorem.

\begin{thmbox}{13.3}
\end{thmbox}

If X is in L2(P) , then E[X\vert\mathcal{G}] is the MMSE estimator of X that achieves the

minimum

min

Z

Z- X = E[X\vert\mathcal{G}]- X

with the minimum taken over all \mathcal{G}-measurable random variables in L2(P) .

The example of nonlinear MMSE estimator in Example 12.4.3 is indeed an

example of conditional expectation in the case of a finite sample space.

Remark 13.1 When \mathcal{G}is generated by a partition of \Omega consisting of k disjoint

events, then E[X\vert\mathcal{G}] reduces to the conditional expectation defined in the previous

section.

We first prove that the condition in (13.3) uniquely defines the conditional

expectation, up to a null set. That is, if there are two solutions to (13.3), then they

must be equal almost surely.

\begin{thmbox}{13.4 (Uniqueness of Conditional Expectation)}
\end{thmbox}

If. Z1 and. Z2 are\mathcal{G}-measurable random variables that satisfy

\int

B

Z1 dP =

\int

B

Z2 dP

for allB \in\mathcal{G}, thenZ1 = Z2 as.

\textit{Proof} Let A be the event .{ : Z1() \geq Z2()}The event A is \mathcal{G}-measurable

becauseY = Z1 - Z2 is\mathcal{G}-measurable andA = Y- 1([0,\infty )).

Because A is in. \mathcal{G}and the expectation operator is linear, we have

0 =

\int

A

Z1 - Z2 dP.

However, Z1() - Z2() \geq 0 for in A. The above equation can hold only when

Z1() = Z2() almost everywhere on A, meaning that there is an event E1 \subset A

withP(E 1) = 0 such thatZ1() = Z2() for all. \in A \ E1.

Likewise, consider the event B ={ : Z1() \leq Z2()}We get Z1() =

Z2() almost everywhere on B. There is an event E2 \subset B with P(E 2) = 0 and

Z1() = Z2() for all. \in B \ E2.

We complete the proof by noting that A \cup B = \Omega andP(E 1 \cup E2) = 0. \hfill $\square$

Theorem 13.4 says that if Y is a random variable that satisfies (13.3), any other

random variable. Y' that satisfies (13.3) is equal to Y with probability 1. On the other

hand, if Y is a random variable that satisfies (13.3), then any random variable that

equals Y as. is also a solution to (13.3)A solution to (13.3) is called a version of

E[X\vert\mathcal{G}].

We illustrate the abstract definition of conditional expectation by considering the

smallest and the largest possible sub-\sigma-algebras.

\textbf{Example 13.2.1 (Conditioning on the Smallest\sigma-Algebra)} \\

Take \mathcal{G}={ \emptyset,\Omega }A \mathcal{G}-measurable random variable X is a constant function. That is, there is a

constant c such thatX() = c for all. \in \Omega We want to determine the value of cI n (13.3),i fw e

takeB = \Omega ,we get

\int

\Omega

E[X\vert{\emptyset,\Omega }]dP = E[X].

On the other hand, since E[X\vert\mathcal{G}]() is equal to the constant c for all. ,we have

\int

\Omega

E[X\vert{\emptyset,\Omega }]dP =

\int

\Omega

cdP = c.

Therefore, the value of c is equal to E[X]The conditional expectation E[X\vert{\emptyset,\Omega }] is a constant

function with valueE[X].

\textbf{Example 13.2.2 (Conditioning on the Largest\sigma-Algebra)} \\

Consider the case \mathcal{G}= \mathcal{F}Let B be any set in \mathcal{G}, ie., B is any set in \mathcal{F}W ew a n tt ofi n d a \mathcal{G}-

measurable function, which is the same as an \mathcal{F}-measurable function E[X\vert\mathcal{F}] such that (13.3)

holds for all B \in\mathcal{F}In fact, X is such a choice for E[X\vert\mathcal{F}]Hence,E[X\vert\mathcal{F}]= X as.

We defineE[X\vertY] in terms of the \sigma-algebra generated by Y .

\begin{defbox}{13.4}
\end{defbox}

The conditional expectation of a random variable X given random variable Y is

defined as

E[X\vertY] \coloneqqE[X\vert\sigma(Y)].

For a family of random variables .{Yi}i\inI , the conditional expectation of X given

{Yi}i\inI is defined as

E[X\vert{Yi}i\inI ] \coloneqqE[X\vert\sigma({Yi}i\inI )].

In particular, the conditional expectation of X given random variables Y and Z

is defined as

E[X\vertY,Z ] \coloneqqE[X\vert\sigma(Y,Z) ].

Conditional probability is defined in terms of conditional expectation.

In the above definition, \sigma(Y,Z) means the smallest \sigma-algebra under which both

Y and Z are measurable, and \sigma({Yi}i\inI ) denotes the smallest \sigma-algebra such that

all random variables. Yi 's are measurable.

\begin{defbox}{13.5}
\end{defbox}

Let .(\Omega, \mathcal{F},P) be a probability space. The conditional probability of an event A

in. \mathcal{F}given a random variable X is defined as

P(A\vertX) \coloneqqE[1A\vert\sigma(X)].

When the conditional expectation E[X\vertY] as regarded as a function of Y ,w e

write it as E[X\vertY = y]This is a function from R to R, with y as the argument.

However, the more fundamental definition of conditional expectation E[X\vertY] is a

function from. \Omega to. R, which is measurable with respect to \sigma-field\sigma(Y) .
